\documentclass[oneside,12pt]{article}          % ----- article

\input{../../Latex/inputs/loading_packages.tex}			% ----- Packages utiles
\input{../../Latex/inputs/preambule_general.tex}         	% ----- Préambule feuilles de TD
\input{../../Latex/inputs/mymacros.tex}                 % ----- Macros mathématiques
\input{../../Latex/inputs/mise_en_page.tex} 		   	% ----- Macros de mise en page

\fancyhead[L]{Séance \no 4}
\fancyhead[R]{T\up{ale} enseignement scientifique}            % En-tête droit
\fancyhead[C]{Suites numériques}
\fancyhf[CF]{\footnotesize{\etablissement}}                     % En-tête central

\titlespacing{\section}{0pt}{2mm}{0pt}
\titlespacing{\subsection}{0pt}{5mm}{1mm}

%%%%%%%%%%%%%%%%%%%%%%%%%%%%%%%%%%%%%%%%%%%%%%%%%%%%%%%%%%%%%%%%%%%%%%%%%%%%%%%%%%%%%%%%%%%%%%%%%%%%%%%%%%%%%
%%%%%%%%%%%%%%%%%%%%%%%%%%%%%%%%%%%%%%%%%%%%%%%%%%%%%%%%%%%%%%%%%%%%%%%%%%%%%%%%%%%%%%%%%%%%%%%%%%%%%%%%%%%%%
%%%%%%%%%%%%%%%%%%%%%%%%%%%%%%%%%%%%%%%%%%%%%%%%%%%%%%%%%%%%%%%%%%%%%%%%%%%%%%%%%%%%%%%%%%%%%%%%%%%%%%%%%%%%%

\begin{document}

\section{Rappels}

\subsection{Gisement de pétrole}
Un gisement de pétrole a produit 60\,000 barils en l'an 2000 et sa production diminue de 1\,500 barils chaque année. On note $u_n$ le nombre de barils produits lors de l'année 2000 + $n$.
\begin{enumerate}[align = left]
	\item Quel est le nombre de barils produits lors des années 2021, 2022 et 2023 ?
	\item Exprimer $u_{n+1}$ en fonction de $u_n$
	\item Quelle formule peut-on utiliser pour calculer directement la $n$-ième terme de la suite ?
	\item [\underline{Bonus} :] Le gisement mettra la clef sous la porte s'il produit moins de 20\,000 barils par an. En quelle année peut-on prévoir sa fermeture ?
\end{enumerate}

\begin{tcolorbox}[
    enhanced , arc = 3pt , top = 4mm ,
    fonttitle = \bfseries , title = Je retiens , boxed title style = {arc = 3pt} ,
    coltitle = black , colbacktitle = white ,
    attach boxed title to top left = {xshift = 5mm , yshift = -4mm} ,
	]
	\restoreskip
	
	Si l'on rajoute (ou soustrait) toujours la même quantité, on parle de suite ..............................

	Pour calculer rapidement ses termes, on utilise la formule : $u_n =$ .........................
\end{tcolorbox}


\subsection{Compte épargne}
Pour la naissance de leur fille, M et Mme Rakotobe placent 1\,000\euro\, sur un compte bancaire à taux d'intérêts annuels de $5\%$. On note $v_n$ le nombre de barils produits lors de l'année 2000 + $n$.
\begin{enumerate}
	\item Quel sera le solde du compte après 1 an ? 2 ans ? 3 ans ?
	\item Exprimer $v_{n+1}$ en fonction de $v_n$
	\item Quelle formule peut-on utiliser pour calculer directement la $n$-ième terme de la suite ?
\end{enumerate}

\begin{tcolorbox}[
    enhanced , arc = 3pt , top = 4mm ,
    fonttitle = \bfseries , title = Je retiens  , boxed title style = {arc = 3pt} ,
    coltitle = black , colbacktitle = white ,
    attach boxed title to top left = {xshift = 5mm , yshift = -4mm} ,
	]
	\restoreskip
	
	Si l'on multiplie (ou divise) toujours par la même quantité, on parle de suite ..............................

	Pour calculer rapidement ses termes, on utilise la formule : $u_n =$ .........................
\end{tcolorbox}



\section{Une suite particulière}

\begin{minipage}{0.55\textwidth} 
	\subsection{La situation} 
	
	Un élève très pressé de ne pas arriver en retard à son cours d'enseignement scientifique doit monter un escalier avant d'arriver dans sa salle de classe.
	
	\bigskip
	Pour aller plus vite, il décide de monter les marches de l'escalier soit par une par une, soit deux par deux.
	
	\bigskip
	À chaque marche, il peut donc choisir de monter soit une seule nouvelle marche, soit deux marches d'un coup !
	
	\bigskip
	Sur le schéma ci-contre on a représenté deux telles manières de monter un escalier comportant 7 marches :
	
	\bigskip
	Mais de \textbf{combien de manières différentes} cet escalier aurait-il pu être gravi ?
\end{minipage}
\hspace{5mm}
\begin{minipage}{0.45\textwidth}
	\begin{tikzpicture}[>=Stealth]
	\def\step{6}
		\foreach \i in {0,...,\step}{
			\draw (\i,\i) -- ++(0,1) -- ++(1,0) ;
			}
		\draw (0,0) -- (\step+1 , 0) -- (\step+1 , \step+1) ;
	
		%%% Chemin bleu %%%
		\draw[newblue, line width=1.5pt, -Stealth , shorten > = 2pt]
		    (-0.5,0.15) to[out=80,in=140,looseness=1.75] (0.5,1.15);
		
		\draw[newblue, line width=1.5pt, -Stealth , shorten > = 2pt]
		    (0.5,1.15) to[out=80,in=130,looseness=1.5] (1.5,2.15);
		
		\draw[newblue, line width=1.5pt, -Stealth , shorten > = 2pt]
		    (1.5,2.15) to[out=80,in=130,looseness=1.5] (3.5,4.5);
		
		\draw[newblue, line width=1.5pt, -Stealth , shorten > = 2pt]
		    (3.5,4.5) to[out=90,in=130,looseness=1.25] (5.5,6.15);
		
		\draw[newblue, line width=1.5pt, -Stealth , shorten > = 2pt]
		    (5.5,6.15) to[out=80,in=120,looseness=1.75] (6.5,7.5);
		    
		%%% Chemin rouge %%%
		\draw[newred, line width=1.5pt, -Stealth , shorten > = 2pt]
		    (-0.5,0.35) to[out=90,in=120,looseness=2] (1.5,2.5);
		
		\draw[newred, line width=1.5pt, -Stealth , shorten > = 2pt]
		    (1.5,2.5) to[out=70,in=130,looseness=1.75] (2.5,3.35);
		
		\draw[newred, line width=1.5pt, -Stealth , shorten > = 2pt]
		    (2.5,3.35) to[out=80,in=130,looseness=1.5] (3.5,4.15);
		    
		\draw[newred, line width=1.5pt, -Stealth , shorten > = 2pt]
		    (3.5,4.15) to[out=80,in=130,looseness=1.5] (4.5,5.15);
		    
		\draw[newred, line width=1.5pt, -Stealth , shorten > = 2pt]
		    (4.5,5.15) to[out=80,in=130,looseness=1.5] (6.5,7.15);
		
	\end{tikzpicture}
\end{minipage}



\subsection{Étude de la situation}
On admettra qu'il existe exactement une manière de monter un escalier à $0$ marches : ne pas monter.
\begin{enumerate}[itemsep = 2mm]
	\item Compléter le tableau suivant :
	
	\renewcommand{\arraystretch}{2} \hspace{-12.5mm}
	\begin{tabular}{| P{4cm} | P{1.75cm} | P{1.75cm} | P{1.75cm} | P{1.75cm} | P{1.75cm} | P{1.75cm} | P{1.75cm} |} \hline
		Nombre de marches & $0$ & $1$ & $2$ & $3$ & $4$ & $5$ & $6$ \\ \hline
		Nombre de manières de monter l'escalier & $1$ & & & & & & \\ \hline
	\end{tabular}

	\item Remarquez-vous un lien entre les nombre de la deuxième ligne ?
	
	\item En déduire la réponse à notre situation initiale :
	
	\begin{tcolorbox}[
	    enhanced , arc = 3pt , top = 4mm , hbox ,
	    fonttitle = \bfseries , title = Réponse , boxed title style = {arc = 3pt} ,
	    coltitle = black , colbacktitle = white ,
	    attach boxed title to top left = {xshift = 5mm , yshift = -4mm} ,
		]
		Il y a exactement ........ manières différentes de monter un escalier à 7 marches.
	\end{tcolorbox}

\end{enumerate}



\subsection{La suite de Fibonacci}
Les nombres que l'on a rentré dans le tableau précédent constituent la très célèbre \textbf{suite de Fibonacci}.

\begin{tcolorbox}[
    enhanced , arc = 3pt , top = 4mm ,
    fonttitle = \bfseries , title = Je retiens , boxed title style = {arc = 3pt} ,
    coltitle = black , colbacktitle = white ,
    attach boxed title to top left = {xshift = 5mm , yshift = -4mm} ,
	]
	\restoreskip
	Dans la suite de Fibonacci, pour obtenir chaque nouveau terme il faut ............................................. \\[10pt]
.............................................................................................................................................................
\end{tcolorbox}



\textbf{\underline{Question} :} Combien de manière il y a-t-il de monter un escalier à 100 marches ? Autrement dit, combien vaut le centième terme de la suite de Fibonacci ?


\begin{tcolorbox}[
    enhanced , arc = 3pt , top = 4mm ,
    fonttitle = \bfseries , title = Réponse par algorithme  , boxed title style = {arc = 3pt} ,
    coltitle = black , colbacktitle = white ,
    attach boxed title to top left = {xshift = 5mm , yshift = -4mm} ,
	]
	\restoreskip
	
	\begin{enumerate}
		\item Choisissez deux nombres quelconques $a$ et $b$ puis exécutez plusieurs fois les instructions suivantes :	
			\begin{codebox}
b = a + b
a = b - a		
			\end{codebox}
		
		\item Comment décrire l'effet de ces instructions sur les nombres $a$ et $b$ ?
		\item En déduire un algorithme qui permet de calculer le centième terme de la suite de Fibonacci.
	\end{enumerate}		
	



	
\end{tcolorbox}




\end{document}